Este trabajo estudia el seguimiento de trayectorias de un cuadricóptero mediante control clásico y control neuronal por imitación. El objetivo no es demostrar que una red neuronal supera siempre a un controlador PID optimizado, sino analizar en qué condiciones una política aprendida puede generalizar mejor que controladores ajustados para familias concretas de trayectorias.

Como parte del alcance se desarrolla un simulador propio de seis grados de libertad. La decisión responde a la necesidad de controlar las hipótesis físicas, los sistemas de referencia, la telemetría y la trazabilidad de los experimentos. Frente a simuladores existentes, normalmente más cerrados, más extensos o menos adaptados a un estudio acotado de control, el simulador desarrollado permite ligar directamente modelo físico, escenarios, datasets y métricas.

La memoria presenta el modelo dinámico, la estrategia de control clásico, los controladores neuronales entrenados por imitación, la generación de escenarios y el protocolo experimental usado para comparar controladores en trayectorias nominales, transferencia entre familias y condiciones fuera de distribución.
